\chapter{Bilangan Kompleks}\label{ch:g12:complex}

\omterm{def:g10:algebra:equation}{Persamaan} $x^2 = -1$ tak punya penyelesaian real. Memperbesar $\R$ dengan satu
bilangan baru $\iu$ yang kuadratnya $-1$ menghasilkan lapangan $\C$ berisi \omterm{def:g12:complex:def}{bilangan
kompleks}, dan di dalamnya \emph{setiap} \omterm{def:g10:algebra:equation}{persamaan} suku banyak punya penyelesaian ---
lagi pula aritmetikanya menyandikan geometri bidang: geseran, putaran
dan penskalaan menjadi penjumlahan dan perkalian.

\section{Bentuk aljabar}

\begin{definition}[Bilangan kompleks]\label{def:g12:complex:def}
Himpunan \emph{bilangan kompleks}\index{bilangan kompleks} adalah
\[
\C = \{a + \iu b : a, b \in \R\},
\]
dengan $\iu$ sebuah lambang yang tunduk pada satu kaidah $\iu^2 = -1$; penjumlahan
dan perkaliannya dikerjakan seperti pada suku banyak real dalam $\iu$, dengan
menyusutkan $\iu^2$ menjadi $-1$. \omterm{def:g10:numbers:sets}{Bilangan real} $a$ dan $b$ disebut \emph{bagian
real}\index{bagian real} $\Rea(z)$ dan \emph{bagian imajiner}\index{bagian imajiner}
$\Ima(z)$ dari $z = a + \iu b$; ungkapan ini disebut \emph{bentuk
aljabar}\index{bentuk aljabar} dari $z$, dan bentuk itu tunggal:
$a + \iu b = a' + \iu b'$ jika dan hanya jika $a = a'$ dan $b = b'$.
\end{definition}

\begin{example}
$(2 + 3\iu)(1 - \iu) = 2 - 2\iu + 3\iu - 3\iu^2 = 5 + \iu$.
\end{example}

\begin{definition}[Konjugat]\label{def:g12:complex:conjugate}
\emph{Konjugat}\index{konjugat} dari $z = a + \iu b$ adalah
$\conj{z} = a - \iu b$.
\end{definition}

\begin{proposition}[Sifat pengonjugatan]\label{prop:g12:complex:conjugate}
Untuk semua $z, w \in \C$:
\[
\conj{z + w} = \conj z + \conj w, \quad
\conj{zw} = \conj z\,\conj w, \quad
\conj{\conj z} = z, \quad
z + \conj z = 2\Rea(z), \quad
z - \conj z = 2\iu\,\Ima(z),
\]
dan $z\conj z = a^2 + b^2 \geq 0$ untuk $z = a + \iu b$. Lagi pula $z \in \R$
jika dan hanya jika $z = \conj z$.
\end{proposition}

\begin{proof}
Semuanya perhitungan langsung dari \omterm{def:g12:complex:def}{bentuk aljabarnya}; misalnya
$z\conj z = (a + \iu b)(a - \iu b) = a^2 - (\iu b)^2 = a^2 + b^2$.
\end{proof}

\begin{method}[Membagi bilangan kompleks]\label{met:g12:complex:division}
Untuk menulis hasil bagi dalam \omterm{def:g12:complex:def}{bentuk aljabar}, kalikanlah pembilang dan penyebutnya
dengan \omterm{def:g12:complex:conjugate}{konjugat} penyebutnya:
\[
\frac{1}{2 + \iu} = \frac{2 - \iu}{(2+\iu)(2-\iu)} = \frac{2 - \iu}{5}
= \frac25 - \frac15\iu .
\]
Karena itu setiap \omterm{def:g12:complex:def}{bilangan kompleks} tak nol punya balikan: $\C$ adalah sebuah
\emph{lapangan}.
\end{method}

\section{Bidang kompleks, modulus dan argumen}

\begin{definition}[Afiks, modulus]\label{def:g12:complex:modulus}
Pada bidang berkoordinat \omterm{def:g10:coordgeom:system}{ortonormal}, titik $M(a, b)$ disebut
\emph{bayangan} dari $z = a + \iu b$, dan $z$ disebut \emph{afiks}\index{afiks} $M$.
\emph{Modulus}\index{modulus} dari $z$ adalah jarak $M$ ke titik asalnya:
\[
\abs{z} = \sqrt{a^2 + b^2} = \sqrt{z\conj z}.
\]
\end{definition}

\begin{omfigure}
\begin{tikzpicture}[scale=1.5]
  \draw[->] (-0.7,0) -- (2.9,0) node[right] {$x$};
  \draw[->] (0,-1.9) -- (0,2.1) node[above] {$y$};
  \draw[dashed] (2,1.5) -- (2,0) node[below right] {$a$};
  \draw[dashed] (2,1.5) -- (0,1.5) node[left] {$b$};
  \draw[dashed] (2,-1.5) -- (2,0);
  \draw[-Stealth,omDef,thick] (0,0) -- (2,1.5)
    node[midway, above left, font=\small] {$\abs z$};
  \draw[-Stealth,omThm,thick] (0,0) -- (2,-1.5);
  \fill (2,1.5) circle (1pt) node[above right] {$M(z = a + \iu b)$};
  \fill (2,-1.5) circle (1pt) node[below right] {$\conj z = a - \iu b$};
\end{tikzpicture}

{\small Titik $M(a,b)$ berafiks $z$: modulusnya $\abs z$ adalah
panjang $OM$, dan \omterm{def:g12:complex:conjugate}{konjugatnya} $\conj z$ adalah pencerminan $z$ pada
sumbu realnya.}
\end{omfigure}

\begin{proposition}[Sifat modulus]\label{prop:g12:complex:modulus}
Untuk $z, w \in \C$:
\[
\abs{zw} = \abs z\,\abs w, \qquad
\abs{\tfrac zw} = \tfrac{\abs z}{\abs w}\ (w \neq 0), \qquad
\abs{\conj z} = \abs z, \qquad
\abs{z + w} \leq \abs{z} + \abs{w} \ \text{(ketaksamaan segitiga)}.
\]
Lagi pula $\abs{z_B - z_A}$ adalah jarak antara titik yang berafiks
$z_A$ dan $z_B$.
\end{proposition}

\begin{proof}
$\abs{zw}^2 = zw\,\conj{zw} = z\conj z\, w \conj w = \abs z^2 \abs w^2$
memberi yang pertama; yang kedua dan ketiga serupa. Untuk ketaksamaan
segitiganya:
\[
\abs{z+w}^2 = (z+w)(\conj z + \conj w)
= \abs z^2 + \abs w^2 + 2\Rea(z \conj w)
\leq \abs z^2 + \abs w^2 + 2\abs{z\conj w}
= \bigl(\abs z + \abs w\bigr)^2,
\]
dengan memakai $\Rea(u) \leq \abs u$. Pernyataan jaraknya adalah rumus
Pythagoras yang diterapkan pada \omterm{def:g10:coordgeom:system}{koordinat} $z_B - z_A$.
\end{proof}

\begin{definition}[Argumen, bentuk trigonometri dan bentuk eksponen]\label{def:g12:complex:argument}
Misalkan $z \neq 0$. Suatu \emph{argumen}\index{argumen} dari $z$, yang ditulis
$\Arg(z)$, adalah sebarang ukuran $\theta$ bagi sudut dari sumbu real positif
ke sinar $OM$; ukuran itu tertentu sampai pada $2k\pi$. Dengan menulis $r = \abs z$,
\[
z = r(\cos\theta + \iu\sin\theta)
\qquad\text{(bentuk trigonometri)}.
\]
Dengan memperkenalkan lambang\index{bentuk eksponen}
\[
\eu^{\iu\theta} = \cos\theta + \iu\sin\theta,
\]
ini menjadi \emph{bentuk eksponen} $z = r\,\eu^{\iu\theta}$.
\end{definition}

\begin{omfigure}
\begin{tikzpicture}[scale=1.5]
  \draw[->] (-0.7,0) -- (2.4,0) node[right] {$x$};
  \draw[->] (0,-0.5) -- (0,1.9) node[above] {$y$};
  \draw[gray, dashed] (1.9,0) arc (0:75:1.9);
  \draw[-Stealth,omDef,thick] (0,0) -- (40:1.9)
    node[pos=0.62, below right=1pt, font=\small] {$r = \abs z$};
  \fill (40:1.9) circle (1pt) node[above right] {$z = r\,\eu^{\iu\theta}$};
  \draw[->,omThm,thick] (0.45,0) arc (0:40:0.45);
  \node[omThm] at (20:0.65) {$\theta$};
\end{tikzpicture}

{\small \omterm{def:g12:complex:argument}{Bentuk eksponen}: $z$ ditentukan oleh jaraknya $r$ ke
titik asal dan sudutnya $\theta$ dari sumbu real positif.}
\end{omfigure}

\begin{theorem}\label{thm:g12:complex:expform}
Untuk semua $\theta, \theta' \in \R$:
\[
\eu^{\iu\theta}\,\eu^{\iu\theta'} = \eu^{\iu(\theta + \theta')},
\qquad
\conj{\eu^{\iu\theta}} = \eu^{-\iu\theta},
\qquad
\abs{\eu^{\iu\theta}} = 1 .
\]
Akibatnya, untuk $z = r\eu^{\iu\theta}$ dan $z' = r'\eu^{\iu\theta'}$
yang tak nol:
\[
zz' = rr'\,\eu^{\iu(\theta+\theta')},
\qquad
\frac{z}{z'} = \frac{r}{r'}\,\eu^{\iu(\theta - \theta')} :
\]
\emph{modulusnya mengalikan, \omterm{def:g12:complex:argument}{argumennya} menjumlahkan.}
\end{theorem}

\begin{proof}
Identitas pertamanya \emph{adalah} sepasang rumus penjumlahan
(\cref{prop:g12:trigo:addition}):
\[
(\cos\theta + \iu\sin\theta)(\cos\theta' + \iu\sin\theta')
= (\cos\theta\cos\theta' - \sin\theta\sin\theta')
+ \iu(\sin\theta\cos\theta' + \cos\theta\sin\theta').
\]
Sisanya menyusul lewat perhitungan langsung.
\end{proof}

\begin{corollary}[Rumus De Moivre]\label{cor:g12:complex:demoivre}
\index{rumus De Moivre}
Untuk $\theta \in \R$ dan $n \in \Z$:
$\bigl(\cos\theta + \iu\sin\theta\bigr)^n = \cos n\theta + \iu \sin n\theta$.
\end{corollary}

\begin{proof}
Induksi pada $n \geq 0$ dengan memakai \cref{thm:g12:complex:expform}; lalu
perluas ke $n < 0$ dengan mengambil balikannya.
\end{proof}

\begin{method}[Berpindah antarbentuk]\label{met:g12:complex:forms}
Dari aljabar ke eksponen: hitunglah $r = \abs z$, lalu carilah $\theta$ dengan
$\cos\theta = \frac ar$, $\sin\theta = \frac br$ (tentukan dahulu kuadran yang benar
sebelum memanggil \omterm{def:g12:trigo:cossin}{kosinus} balikannya). Dari eksponen ke aljabar: jabarkan
$r\cos\theta + \iu\, r\sin\theta$. Pakailah \omterm{def:g12:complex:def}{bentuk aljabar} untuk penjumlahan, dan
\omterm{def:g12:complex:argument}{bentuk eksponen} untuk perkalian, pemangkatan dan pembagian.
\end{method}

\begin{example}
$z = 1 + \iu$: $\abs z = \sqrt2$ dan $\cos\theta = \sin\theta =
\frac{1}{\sqrt2}$, jadi $\theta = \frac\pi4$ dan $z = \sqrt2\,\eu^{\iu\pi/4}$.
Karena itu $z^{20} = 2^{10}\,\eu^{5\iu\pi} = -1024$.
\end{example}

\section{Persamaan derajat dua dan akar satuan}

\begin{theorem}[Persamaan kuadrat berkoefisien real]\label{thm:g12:complex:quadratic}
Misalkan $a, b, c \in \R$, $a \neq 0$, dan $\Delta = b^2 - 4ac$. \omterm{def:g10:algebra:equation}{Persamaan}
$az^2 + bz + c = 0$ mempunyai penyelesaian di $\C$:
\[
\begin{aligned}
&\Delta > 0:\ z = \frac{-b \pm \sqrt{\Delta}}{2a} \ (\text{dua yang real}); \qquad
\Delta = 0:\ z = \frac{-b}{2a} \ (\text{kembar}); \\
&\Delta < 0:\ z = \frac{-b \pm \iu\sqrt{-\Delta}}{2a} \ (\text{dua yang berkonjugat}).
\end{aligned}
\]
\end{theorem}

\begin{proof}
Lengkapkanlah kuadratnya:
$az^2 + bz + c = a\left(\left(z + \frac{b}{2a}\right)^2 -
\frac{\Delta}{4a^2}\right)$. Jika $\Delta < 0$, maka
$\frac{\Delta}{4a^2} = \left(\iu\frac{\sqrt{-\Delta}}{2a}\right)^2$, dan
selisih kuadratnya terfaktorkan seperti biasa.
\end{proof}

\begin{definition}[Akar satuan]\label{def:g12:complex:rootsofunity}
Untuk $n \geq 1$, yang disebut \emph{akar satuan ke-$n$}\index{akar satuan} adalah
penyelesaian $z^n = 1$.
\end{definition}

\begin{theorem}\label{thm:g12:complex:rootsofunity}
\omterm{def:g10:algebra:equation}{Persamaan} $z^n = 1$ mempunyai tepat $n$ penyelesaian:
\[
\omega_k = \eu^{2\iu k\pi/n}, \qquad k = 0, 1, \dots, n-1 .
\]
Bayangannya adalah titik sudut segi-$n$ beraturan yang termuat dalam lingkaran
satuan, dan untuk $n \geq 2$ jumlahnya $0$.
\end{theorem}

\begin{proof}
Setiap $\omega_k$ memenuhi $\omega_k^n = \eu^{2\iu k\pi} = 1$, dan semuanya
berbeda sepasang demi sepasang sebab \omterm{def:g12:complex:argument}{argumennya} terletak di $\intco{0}{2\pi}$.
Sebaliknya, jika $z^n = 1$ maka $\abs z^n = 1$, sehingga $\abs z = 1$ (\omterm{def:g10:numbers:sets}{bilangan
real} positif) dan $z = \eu^{\iu\theta}$ dengan $n\theta \equiv 0 \pmod{2\pi}$,
\emph{yaitu} $\theta = \frac{2k\pi}{n}$; menyusutkan $k$ modulo $n$ mendarat di
antara $\omega_k$ itu. Titik sudutnya berjarak sudut sama $\frac{2\pi}{n}$:
itulah segi-$n$ beraturan. Akhirnya, dengan $\omega = \omega_1$,
berlaku $\omega \neq 1$ dan jumlah deret ukurnya memberi
\[
\sum_{k=0}^{n-1}\omega_k = \sum_{k=0}^{n-1}\omega^k
= \frac{\omega^n - 1}{\omega - 1} = 0 . \qedhere
\]
\end{proof}

\begin{omfigure}
\begin{tikzpicture}[scale=1.7]
  \draw[->] (-1.3,0) -- (1.45,0) node[right] {$x$};
  \draw[->] (0,-1.3) -- (0,1.4) node[above] {$y$};
  \draw[gray] (0,0) circle (1);
  \draw[omDef] (0:1) -- (72:1) -- (144:1) -- (216:1) -- (288:1) -- cycle;
  \draw[->,omThm,thick] (0.3,0) arc (0:72:0.3);
  \node[omThm, font=\small] at (36:0.52) {$\frac{2\pi}{5}$};
  \fill (0:1) circle (1pt) node[below right] {$1$};
  \fill (72:1) circle (1pt) node[above right] {$\omega$};
  \fill (144:1) circle (1pt) node[above left] {$\omega^2$};
  \fill (216:1) circle (1pt) node[below left] {$\omega^3$};
  \fill (288:1) circle (1pt) node[below right] {$\omega^4$};
\end{tikzpicture}

{\small \omterm{def:g12:complex:rootsofunity}{Akar satuan} kelima, $\omega = \eu^{2\iu\pi/5}$: yaitu
titik sudut segi lima beraturan yang termuat dalam lingkaran satuan
(\cref{exo:g12:complex:10}).}
\end{omfigure}

\begin{example}
\omterm{def:g11:quad:discriminant}{Akar} pangkat tiga dari satuan adalah $1$, $j = \eu^{2\iu\pi/3} =
-\frac12 + \iu\frac{\sqrt3}{2}$ dan $\conj j = j^2$, yang memenuhi
$1 + j + j^2 = 0$.
\end{example}

\section{Bilangan kompleks dan geometri bidang}

\begin{proposition}[Kamus geometris]\label{prop:g12:complex:geometry}
Misalkan $A, B, C, D$ titik yang berafiks $a, b, c, d$ ($a \neq b$,
$c \neq d$).
\begin{enumerate}
  \item Geseran oleh \omterm{def:g11:vect:vector}{vektor} berafiks $t$ adalah $z \mapsto z + t$.
  \item Putaran berpusat $\omega$ (afiksnya) dan bersudut $\theta$ adalah
        $z \mapsto \omega + \eu^{\iu\theta}(z - \omega)$.
  \item Penskalaan (homoteti) berpusat $\omega$ dan berrasio
        $k \in \R^*$ adalah $z \mapsto \omega + k(z - \omega)$.
  \item $\dfrac{d - c}{b - a}$ bermodulus $\dfrac{CD}{AB}$ dan berargumen
        sudut antara \omterm{def:g11:vect:vector}{vektor} $\vect{AB}$ dan $\vect{CD}$. Khususnya,
        $AB \perp CD$ jika dan hanya jika $\frac{d-c}{b-a}$ murni
        imajiner, dan $A, B, C$ segaris ($C \ne A$, $C\ne B$) jika dan
        hanya jika $\frac{c-a}{b-a} \in \R$.
\end{enumerate}
\end{proposition}

\begin{proof}
(1) dan (3) hanyalah rumus \omterm{def:g10:coordgeom:system}{koordinatnya} yang dinyatakan ulang. Untuk (2), pemetaan
$w \mapsto \eu^{\iu\theta} w$ mengalikan modulusnya dengan $1$ dan menambahkan
$\theta$ pada \omterm{def:g12:complex:argument}{argumennya}: itulah putaran bersudut $\theta$ terhadap titik asal;
mengonjugatkannya dengan geseran yang mengirim $\omega$ ke $0$ memberi pusat yang
umum. Untuk (4), \omterm{def:g12:complex:modulus}{modulus} dan \omterm{def:g12:complex:argument}{argumen} hasil baginya saling mengurangkan
(\cref{thm:g12:complex:expform}).
\end{proof}

\section{Latihan}

\begin{exercise}[$\star$]\label{exo:g12:complex:1}
Tulislah dalam \omterm{def:g12:complex:def}{bentuk aljabar}:
$(3 - 2\iu)(1 + 4\iu)$, \quad $\dfrac{2 + \iu}{1 - \iu}$, \quad
$\iu^{2027}$.
\end{exercise}

\begin{exercise}[$\star$]\label{exo:g12:complex:2}
Selesaikan di $\C$: (a) $z^2 - 4z + 13 = 0$; (b) $z^2 + z + 1 = 0$. Periksalah
pada tiap kasusnya bahwa kedua penyelesaiannya berkonjugat lalu hitung hasil kalinya.
\end{exercise}

\begin{exercise}[$\star$]\label{exo:g12:complex:3}
Nyatakan dalam \omterm{def:g12:complex:argument}{bentuk eksponen}: $z_1 = -1 + \iu$, $z_2 = \sqrt3 - \iu$, lalu hitung
$\dfrac{z_1}{z_2}$ dalam \omterm{def:g12:complex:argument}{bentuk eksponen} lalu dalam \omterm{def:g12:complex:def}{bentuk aljabar}. Simpulkan nilai
eksak $\cos\dfrac{11\pi}{12}$.
\end{exercise}

\begin{exercise}[$\star\star$]\label{exo:g12:complex:4}
Lukiskan secara geometris himpunan titik $M$ berafiks $z$ sedemikian sehingga:
(a) $\abs{z - 2} = \abs{z + \iu}$;\quad
(b) $\abs{z - 1 - \iu} = 3$;\quad
(c) $\dfrac{z - 1}{z + 1}$ murni imajiner.
\end{exercise}

\begin{exercise}[$\star\star$]\label{exo:g12:complex:5}
Dengan rumus De Moivre dan teorema binomial, nyatakan $\cos 3\theta$
sebagai suku banyak dalam $\cos\theta$, dan $\sin 3\theta$ dengan
$\sin\theta$.
\end{exercise}

\begin{exercise}[$\star\star$]\label{exo:g12:complex:6}
\emph{(Pelinearan.)} Dengan memakai
$\cos\theta = \frac{\eu^{\iu\theta} + \eu^{-\iu\theta}}{2}$, linearkanlah
$\cos^3\theta$ (tulislah sebagai gabungan $\cos k\theta$), lalu simpulkan
$\displaystyle\int_0^{\pi/2}\cos^3\theta\,\dd\theta$.
\end{exercise}

\begin{exercise}[$\star\star$]\label{exo:g12:complex:7}
Misalkan $A$ dan $B$ berafiks $a = 1 + \iu$ dan $b = 3 + 2\iu$. Tentukan kedua
titik $C$ yang membuat segitiga $ABC$ sama sisi, sebagai bayangan $B$
di bawah putaran berpusat $A$ dan bersudut $\pm\frac{\pi}{3}$.
\end{exercise}

\begin{exercise}[$\star\star\star$]\label{exo:g12:complex:8}
Selesaikan $z^4 = -4$ di $\C$ lalu gambarkan penyelesaiannya. Faktorkanlah
$z^4 + 4$ menjadi dua suku banyak kuadrat yang berkoefisien real.
\end{exercise}

\begin{exercise}[$\star\star\star$]\label{exo:g12:complex:9}
Untuk $\theta \in \intoo{0}{2\pi}$ dan $n \in \N$, hitunglah
\[
S = 1 + \cos\theta + \cos 2\theta + \dots + \cos n\theta
\]
dengan menjumlahkan deret ukurnya
$\sum_{k=0}^n \eu^{\iu k\theta}$ lalu mengambil \omterm{def:g12:complex:def}{bagian realnya}. (Jawabannya:
$S = \dfrac{\sin\frac{(n+1)\theta}{2}}{\sin\frac{\theta}{2}}
\cos\frac{n\theta}{2}$.)
\end{exercise}

\begin{exercise}[$\star\star\star$]\label{exo:g12:complex:10}
Misalkan $\omega = \eu^{2\iu\pi/5}$.
\begin{enumerate}
  \item Berilah alasan bahwa $1 + \omega + \omega^2 + \omega^3 + \omega^4 = 0$.
  \item Misalkan $u = \omega + \omega^4$ dan $v = \omega^2 + \omega^3$. Tunjukkan
        bahwa $u + v = -1$ dan $uv = -1$, lalu simpulkan $u = \frac{-1+\sqrt5}{2}$.
  \item Simpulkan bahwa $\cos\dfrac{2\pi}{5} = \dfrac{\sqrt5 - 1}{4}$.
\end{enumerate}
\end{exercise}

\section{Soal: Segi lima di cermin satuan}

\begin{problem}[{Soal akhir pekan --- akar satuan kelima
menghitung $\cos 72^\circ$ secara eksak, nisbah emas menandatangani
hasilnya, dan perkalian ternyata sebuah putaran}]
\label{pb:g12:complex:1}
Euklides sanggup melukis segi lima yang beraturan, tetapi alasan
\emph{mengapa} bangun itu tunduk pada penggaris dan jangka --- sedangkan
sudut $20^\circ$ yang sederhana tidak --- tersembunyi selama dua ribu
tahun. Alasannya bersembunyi di bab ini: kelima \omterm{def:g12:complex:rootsofunity}{akar satuan} kelima
(\cref{thm:g12:complex:rootsofunity}) berjumlah nol, dan satu
baris aljabar itu memeras $\cos 72^\circ$ keluar dari sebuah \omterm{def:g10:algebra:equation}{persamaan}
kuadrat --- hanya berisi akar kuadrat, jadi terlukiskan, dan tentu saja
ditandatangani oleh nisbah emas. Di sekeliling permata ini: kelancaran, jalan
pintas de Moivre menuju trigonometri, dan perkalian yang tersingkap sebagai
geometri.

\medskip
\noindent\textbf{Bagian I --- Kelancaran.}
\begin{enumerate}
  \item Hitunglah $(2 + \iu)(3 - \iu)$, \
        $\dfrac{1 + \iu}{1 - \iu}$
        (\cref{met:g12:complex:division}), dan
        $\abs{3 + 4\iu}$.
  \item Nyatakan dalam \omterm{def:g12:complex:argument}{bentuk eksponen}: $1 + \iu$; \ $-2$; \
        $1 - \iu\sqrt3$.
  \item Hitunglah $(1 + \iu)^8$ dengan memakai \omterm{def:g12:complex:argument}{bentuk eksponennya}.
  \item Selesaikan $z^2 = \iu$, lalu $z^2 - 2z + 5 = 0$
        (\cref{thm:g12:complex:quadratic}).
  \item Kamus geometrisnya
        (\cref{prop:g12:complex:geometry}): lukiskan pemetaan
        $z \mapsto \iu z$, dan himpunan
        $\abs{z - (1 + \iu)} = 2$.
\end{enumerate}

\medskip
\noindent\textbf{Bagian II --- Lima \omterm{def:g11:quad:discriminant}{akar}, satu segi lima.}
Misalkan $\omega = \eu^{2\iu\pi/5}$ lalu tinjau
$1, \omega, \omega^2, \omega^3, \omega^4$.
\begin{enumerate}[resume]
  \item Periksalah bahwa itu tepat penyelesaian $z^5 = 1$,
        lalu lukiskan bangun yang digambarnya pada bidang.
  \item Buktikan bahwa jumlahnya nol. (Lewat deret ukur
        --- atau kalikan jumlahnya dengan $1 - \omega$.)
  \item Ambillah \omterm{def:g12:complex:def}{bagian realnya} lalu simpulkan
        \[
        1 + 2\cos\frac{2\pi}{5} + 2\cos\frac{4\pi}{5} = 0 .
        \]
  \item Ambil $c = \cos\frac{2\pi}{5}$ lalu pakailah rumus sudut
        rangkap bagi $\cos\frac{4\pi}{5}$ untuk mengubah pertanyaan 8
        menjadi \omterm{def:g10:algebra:equation}{persamaan} kuadrat dalam $c$; selesaikan lalu
        simpulkan
        \[
        \cos 72^\circ = \frac{\sqrt5 - 1}{4} .
        \]
  \item Periksalah secara numeris, lalu biarkan nisbah emas menandatangani
        karyanya: tunjukkan $c = \frac{1}{2\varphi}$ dengan
        $\varphi = \frac{1 + \sqrt5}{2}$
        (\cref{pb:g10:algebra:1}) --- lalu nyatakan gema
        geometrisnya yang masyhur: pada segi lima beraturan, diagonalnya
        adalah $\varphi$ kali sisinya.
\end{enumerate}

\medskip
\noindent\textbf{Bagian III --- Jalan pintas de Moivre.}
\begin{enumerate}[resume]
  \item Vonis pelukisannya: $\cos 72^\circ$ hanya memerlukan
        akar kuadrat (pertanyaan 9), sehingga segi lima itu
        terlukiskan dengan penggaris dan jangka; sedangkan $\cos 20^\circ$
        memenuhi \omterm{def:g10:algebra:equation}{persamaan} pangkat tiga yang tak terpecahkan
        (\cref{pb:g12:trigo:1}), sehingga sudut $60^\circ$ tak dapat
        dibagi tiga. Nyatakan kaidah yang mulai muncul itu (akar
        kuadrat terbangun, akar pangkat tiga tidak), yang teori lengkapnya ---
        Gauss, pada usia 19 tahun, dan segi tujuh belas --- hidup dalam
        jilid universitas.
  \item Selesaikan $z^4 = -4$ lewat \omterm{def:g12:complex:argument}{bentuk eksponen} lalu perolehlah kembali
        pemfaktoran pada \cref{exo:g12:complex:8},
        $z^4 + 4 = (z^2 - 2z + 2)(z^2 + 2z + 2)$, dengan memasangkan
        \omterm{def:g11:quad:discriminant}{akar} yang berkonjugat.
  \item Jabarkan $(\cos\theta + \iu\sin\theta)^3$ dengan
        teorema binomial dan de Moivre, lalu turunkan kembali
        $\cos 3\theta = 4\cos^3\theta - 3\cos\theta$ dalam tiga
        baris --- bandingkan dengan jalur rumus penjumlahan pada
        \cref{pb:g12:trigo:1}.
  \item Periksa silang dengan \cref{exo:g12:complex:9}: untuk
        $\theta = \frac{2\pi}{5}$ dan $n = 4$, berapakah nilai
        jumlah $1 + \cos\theta + \dots + \cos 4\theta$, dan
        mengapa itu selaras dengan pertanyaan 8?
  \item Hitunglah ketiga akar pangkat tiga dari $8\iu$ dalam \omterm{def:g12:complex:def}{bentuk
        aljabar}.
\end{enumerate}

\medskip
\noindent\textbf{Bagian IV --- Perkalian adalah geometri.}
\begin{enumerate}[resume]
  \item Lukiskan selengkapnya pemetaan $z \mapsto (1 + \iu)z$:
        sebesar sudut apa ia memutar, dengan faktor berapa ia
        menskalakan, dan apa yang dilakukannya pada persegi satuan?
  \item Buktikan $\abs{z_1 z_2} = \abs{z_1}\,\abs{z_2}$ dengan
        \omterm{def:g12:complex:conjugate}{konjugat} (\cref{prop:g12:complex:modulus}), lalu
        hitunglah $(1+\iu)^n$ untuk $n = 1, 2, 3, 4$ lalu lukiskan
        ulir yang digambar pangkatnya.
  \item Bilangan $\mathrm{j} = \eu^{2\iu\pi/3}$: tunjukkan
        $1 + \mathrm{j} + \mathrm{j}^2 = 0$, lalu periksalah pada segitiga
        $\left(1, \mathrm{j}, \mathrm{j}^2\right)$ kaidah klasiknya:
        $a + \mathrm{j} b + \mathrm{j}^2 c = 0$ berlaku bagi segitiga sama
        sisi $abc$ (pada salah satu arah putarannya).
  \item Teorema dasar aljabar (yang diterima saja:
        d'Alembert--Gauss): setiap suku banyak terfaktorkan sepenuhnya
        atas $\C$. Simpulkan akibat wajarnya di dunia nyata --- setiap
        suku banyak real terfaktorkan menjadi bagian linear dan kuadrat
        yang real --- lalu tunjuklah pertanyaan 12 sebagai contohnya.
  \item Penutup --- potret $\C$, masing-masing satu kalimat: bilangan
        menjadi titik; perkalian menjadi
        putaran-dan-penskalaan; \omterm{def:g12:complex:rootsofunity}{akar satuan} menjadi segi banyak
        beraturan; de Moivre mengubah pemangkatan menjadi identitas
        trigonometri; dan ketertutupan aljabarnya menjamin bahwa tak ada
        \omterm{def:g10:algebra:equation}{persamaan} yang memerlukan dunia yang lebih besar. Koda: dua
        bilangan masyhur mana yang bertemu pada pertanyaan 10?

\end{enumerate}
\end{problem}
